% Chapter 5: Analysis — ablation, cold-start, hard-residual, efficiency, error analysis
\selectlanguage{greek}


\chapter{Ανάλυση}
\label{chap:analysis}


\section{Μελέτη Εκτομής}
\label{sec:ablation}

Σαρώνουμε ένα πλέγμα τιμών $\tau_1 \in \{0{,}6, 0{,}7, 0{,}8, 0{,}9, 0{,}95\}$ και
$\tau_2 \in \{0{,}5, 0{,}6, 0{,}7, 0{,}8, 0{,}9\}$ (25 συνδυασμοί) στο ήδη εκπαιδευμένο
\en{CascadeLP}, χωρίς επανεκπαίδευση των ταξινομητών ανά επίπεδο --- τα όρια
επηρεάζουν μόνο τη δρομολόγηση κατά το συμπερασμό (\S\ref{sec:cascade-routing}) --- και
καταγράφουμε το \en{Macro F1} και το ποσοστό κλήσεων ανά επίπεδο στο σύνολο
επικύρωσης για κάθε συνδυασμό.

Το όριο $\tau_1$ παρουσιάζει μια απότομη ασυνέχεια: για $\tau_1 \in \{0{,}6, 0{,}7\}$
το Επίπεδο~1 χειρίζεται 78{,}34\% των ζευγών (\en{Macro F1} 0{,}9974--0{,}9976), ενώ
για $\tau_1 \geq 0{,}8$ το ποσοστό αυτό καταρρέει απότομα στο 19{,}38\% (και περαιτέρω
στο 17{,}47\% για $\tau_1 \geq 0{,}9$), με ταυτόχρονη \textit{βελτίωση} του συνολικού
\en{Macro F1} στο 0{,}9983--0{,}9987. Αυτό είναι αναμενόμενο δεδομένης της ακραίας
ανισορροπίας κλάσεων στο εμπιστευτικό υποσύνολο του Επιπέδου~1
(\S\ref{sec:cascade-results}): χαλαρώνοντας το όριο επιτρέπουμε σε περισσότερα ζεύγη
να επιλυθούν από έναν ταξινομητή που ουσιαστικά προβλέπει πάντα τη θετική ετικέτα σε
αυτό το καθεστώς, αυξάνοντας τον κίνδυνο σφάλματος στην σπάνια αρνητική κλάση.

Το όριο $\tau_2$ έχει μονότονη επίδραση: αυξάνοντάς το από 0{,}5 σε 0{,}9 (σταθερού
$\tau_1$), το \en{Macro F1} μειώνεται σταδιακά --- π.χ. από 0{,}9987 σε 0{,}9983 για
$\tau_1 = 0{,}8$ --- ενώ το ποσοστό κλήσεων Επιπέδου~3 αυξάνεται από 0\% σε 0{,}66\%.
Αυτό αντικατοπτρίζει άμεσα την αναξιοπιστία του \texttt{\en{EmbeddingClassifier}} στο
υπόλειμμα ζευγών που φτάνει ως εκεί (\en{Macro F1} 0{,}2802, \S\ref{sec:hard-residual}):
κάθε επιπλέον ζεύγος που κλιμακώνεται στο Επίπεδο~3 αντί να επιλυθεί στο Επίπεδο~2
είναι, κατά μέσο όρο, ζημιογόνο για τη συνολική βαθμολογία. Ο καλύτερος συνδυασμός στο
πλέγμα είναι $\tau_1 \in \{0{,}8, 0{,}9, 0{,}95\}$ με $\tau_2 = 0{,}5$
(\en{Macro F1} = 0{,}9987, μηδενικό ποσοστό κλήσεων Επιπέδου~3), ελαφρώς πάνω από τους
προεπιλεγμένους $\tau_1 = 0{,}8$, $\tau_2 = 0{,}7$ (\en{Macro F1} = 0{,}9986). Η
διαφορά είναι μικρή (0{,}0001) και εντός του θορύβου δειγματοληψίας μεγέθους δείγματος
187.603 ζευγών, οπότε δεν αλλάζουμε την προεπιλεγμένη διαμόρφωση --- το βασικό
συμπέρασμα είναι ότι η αλυσίδα είναι ανθεκτική σε ένα ευρύ φάσμα ορίων γύρω από τις
προεπιλογές, με μοναδική εξαίρεση την απότομη ασυνέχεια στο $\tau_1 \approx 0{,}75$.

\subsection{Χαρακτηριστικά \en{Node2Vec} στο Επίπεδο~1}

Εκπαιδεύουμε δύο εκδοχές του \texttt{\en{StructuralClassifier}} στην ίδια διαίρεση
εκπαίδευσης/επικύρωσης: μία στα τέσσερα ευρετικά χαρακτηριστικά (όπως στην
τελική διαμόρφωση, \S\ref{sec:cascade-routing}) και μία στο 68-διάστατο διάνυσμα που
προσαρτά το γινόμενο \en{Hadamard} των ενσωματώσεων \en{Node2Vec} των δύο τελικών
κόμβων (\S\ref{sec:cascade-routing}). Ο πίνακας~\ref{tab:node2vec-ablation} συνοψίζει
την επίδραση.

\begin{table}[h!]
  \centering
  \begin{tabular}{lrr}
    \toprule
    \textbf{Μετρική} & \textbf{Χωρίς \en{Node2Vec}} & \textbf{Με \en{Node2Vec}} \\
    \midrule
    Ποσοστό κλήσεων Επιπέδου~1     & 19{,}38\% & 99{,}30\% \\
    Ακρίβεια (συνολική)            & 0{,}7900  & 0{,}9958  \\
    \en{Macro F1} (συνολική)       & 0{,}7693  & 0{,}9958  \\
    Ακρίβεια (εμπιστευτικό υποσ.)  & 0{,}9985  & 0{,}9970  \\
    \en{Macro F1} (εμπιστευτικό υποσ.) & 0{,}4996 & 0{,}9970 \\
    \bottomrule
  \end{tabular}
  \caption{\en{StructuralClassifier} με και χωρίς χαρακτηριστικά \en{Node2Vec},
  σύνολο επικύρωσης (n=187.603).}
  \label{tab:node2vec-ablation}
\end{table}

Το αποτέλεσμα αντικρούει την αρχική υπόθεση σχεδιασμού (\S\ref{sec:cascade-routing})
ότι το \en{Node2Vec} προσφέρει μόνο οριακή βελτίωση. Η προσάρτηση των ενσωματώσεων
αυξάνει το ποσοστό κλήσεων Επιπέδου~1 από 19{,}38\% σε 99{,}30\% --- σχεδόν κάθε ζεύγος
γίνεται πλέον επιλύσιμο με υψηλή αξιοπιστία στο φθηνότερο επίπεδο --- και, ακόμη
σημαντικότερο, διορθώνει την κατάρρευση του \en{Macro F1} στο εμπιστευτικό υποσύνολο
που παρατηρήθηκε στην \S\ref{sec:cascade-results}: από 0{,}4996 (πρακτικά τυχαία
απόδοση λόγω ανισορροπίας κλάσεων) σε 0{,}9970. Η εξήγηση συνδέεται άμεσα με την
αραιότητα του δομικού γράφου (μέσος βαθμός κόμβου 1{,}30, \S\ref{sec:separability-results}):
τα τέσσερα ευρετικά χαρακτηριστικά βλέπουν μόνο άμεσους κοινούς γείτονες, που
απουσιάζουν από το 99{,}9963\% των ζευγών επικύρωσης (\S\ref{sec:cold-start-results}),
ενώ οι ενσωματώσεις \en{Node2Vec} κωδικοποιούν εγγύτητα μέσω πολυβηματικών τυχαίων
περιπάτων που διασχίζουν το γράφο πέρα από την άμεση γειτονιά, παρέχοντας διακριτικό
σήμα ακριβώς στο καθεστώς όπου τα ευρετικά χαρακτηριστικά είναι εκφυλισμένα.

Παρά το μέγεθος αυτής της βελτίωσης, η τελική διαμόρφωση του \en{CascadeLP} δεν
ενσωματώνει το \en{Node2Vec} (\S\ref{sec:cascade-routing}): η εκπαίδευσή του σε αυτόν
τον γράφο (668.452 κόμβοι) απαίτησε άνω της μίας ώρας υπολογιστικού χρόνου τοίχου
υπό αυστηρό όριο μνήμης, έναντι λίγων δευτερολέπτων για τα ευρετικά
χαρακτηριστικά, και η πλήρης ενσωμάτωσή του στην αλυσίδα θα απαιτούσε νέα
βαθμονόμηση των ορίων $\tau_1, \tau_2$ και επανάληψη όλων των πειραμάτων των
Κεφαλαίων~\ref{chap:experiments}--\ref{chap:analysis}. Το εύρημα καταγράφεται εδώ ως
σαφής κατεύθυνση μελλοντικής εργασίας (Κεφάλαιο~\ref{chap:conclusion}) και όχι ως
αλλαγή στην τρέχουσα διαμόρφωση.


\section{Ανάλυση Εκκίνησης εκ του Μηδενός}
\label{sec:cold-start-results}

Χρησιμοποιώντας τη \texttt{\en{cold\_start\_mask}} (μηδενικοί κοινοί γείτονες μεταξύ
$u$ και $v$) στο σύνολο επικύρωσης, βρίσκουμε ότι το 99{,}9963\% των ζευγών
(187.596 από 187.603, εξαιρουμένων των αυτοβρόχων) είναι εκκίνησης εκ του μηδενός.
Μόλις 7 ζεύγη επικύρωσης έχουν έστω έναν κοινό γείτονα, και τα 7 ανήκουν στην
κατηγορία υψηλού \en{CN} και επιλύονται σωστά στο Επίπεδο~1. Στην πράξη, η
απόδοση ανά επίπεδο στο υποσύνολο εκκίνησης εκ του μηδενός ταυτίζεται αριθμητικά με
την συνολική απόδοση ανά επίπεδο του πίνακα~\ref{tab:cascade-tier-results}: το
Επίπεδο~1 χειρίζεται το 19{,}37\% αυτού του υποσυνόλου με ακρίβεια 0{,}9985 και
\en{Macro F1} 0{,}4996, το Επίπεδο~2 το 80{,}58\% με ακρίβεια 0{,}9991, και το
Επίπεδο~3 το 0{,}05\% με ακρίβεια μόλις 0{,}2857.

Αυτό το εύρημα αναθεωρεί ουσιαστικά τον τρόπο που πρέπει να αντιμετωπίζεται η
έννοια της εκκίνησης εκ του μηδενός σε αυτό το σύνολο δεδομένων. Δεν πρόκειται για
ένα δυσδιάκριτο άκρο της κατανομής, αλλά για τον σχεδόν απόλυτο κανόνα: ο δομικός
γράφος, κατασκευασμένος αποκλειστικά από τις θετικές ακμές εκπαίδευσης
(\S\ref{sec:cascade-routing}), είναι τόσο αραιός (μέσος βαθμός κόμβου 1{,}30· βλ.\
\S\ref{sec:separability-results}) που η ύπαρξη κοινού γείτονα είναι η εξαίρεση, όχι
ο κανόνας. Ένα προηγούμενο κείμενο αυτής της εργασίας ανέφερε ότι «περίπου το 78\%
των νέων θετικών ακμών» εμπλέκουν κόμβο εκκίνησης εκ του μηδενός, βασισμένο σε μια
αναφορά που αποδείχθηκε ανεπιβεβαίωτη και αφαιρέθηκε (\texttt{\en{plan.md}}). Η
πραγματική μετρηθείσα τιμή σε αυτό το σύνολο δεδομένων είναι σχεδόν 100\%, ένα
μέγεθος τάξης διαφορετικό από τον αναφερόμενο αριθμό. Το \en{CascadeLP} δεν
χρειάζεται ξεχωριστό χειρισμό εκκίνησης εκ του μηδενός επειδή το Επίπεδο~1 ήδη
κλιμακώνει αυτόματα κάθε ζεύγος χωρίς δομικό σήμα (\S\ref{sec:cascade-routing}) ---
το γεγονός ότι η απόδοση εκκίνησης εκ του μηδενός ταυτίζεται με τη συνολική απόδοση
επιβεβαιώνει ότι αυτός ο μηχανισμός λειτουργεί όπως σχεδιάστηκε, όχι ότι αγνοεί το
πρόβλημα.


\section{Ανάλυση του Δύσκολου Υπολείμματος}
\label{sec:hard-residual}

Το υπόλειμμα ζευγών που φτάνει στο Επίπεδο~3 του \en{CascadeLP} \textit{και}
χαρακτηρίζεται «δύσκολο» από τον χαρακτηρισμό διαχωρισιμότητας
(\S\ref{sec:forensics-protocol}) αντιπροσωπεύει το υποσύνολο ζευγών για το οποίο
οι φθηνότερες οικογένειες χαρακτηριστικών γνησίως αποτυγχάνουν: 78 ζεύγη στο σύνολο
επικύρωσης (\S\ref{sec:cascade-results}), με ακρίβεια μόλις 0{,}1667.

Πριν από την ποιοτική εξέταση, αξίζει να αποκλειστεί μια μη-σημασιολογική εξήγηση: 37
από τα 78 ζεύγη (47{,}4\%) περιλαμβάνουν έναν κόμβο με πλήρως ελλιπές κείμενο άρθρου
(τιμή \en{NaN} στο \texttt{\en{nodes.tsv}}), οπότε η ενσωμάτωση \en{sentence-transformer}
δεν διαθέτει ουσιαστικό σήμα εισόδου. Ένας μόνο κόμβος (\texttt{\en{id=415071}}, χωρίς
κείμενο) συμμετέχει σε 18 από αυτά τα ζεύγη· σε όλα τα 18 η πραγματική ετικέτα είναι
αρνητική, αλλά ο ταξινομητής προβλέπει θετική ετικέτα στο 55{,}6\% από αυτά ---
ουσιαστικά τυχαία πρόβλεψη λόγω απουσίας δεδομένων, όχι αδυναμία σημασιολογικής
συλλογιστικής. Ωστόσο, η ελλιπής κάλυψη κειμένου \textit{δεν} είναι η κύρια αιτία της
συνολικής κατάρρευσης του Επιπέδου~3: η ακρίβεια στα 41 ζεύγη με πλήρες κείμενο και
στις δύο πλευρές είναι 0{,}0732, χειρότερη από την ακρίβεια 0{,}2703 στα 37 ζεύγη με
ελλιπές κείμενο. Το γνησίως δύσκολο υποσύνολο είναι εκείνο όπου \textit{υπάρχει}
πλούσιο κείμενο και ο ταξινομητής εξακολουθεί να αποτυγχάνει.

Η χειρωνακτική εξέταση δείγματος από αυτά τα 41 ζεύγη πλήρους κειμένου αναδεικνύει
τρία επαναλαμβανόμενα μοτίβα αποτυχίας:

\begin{itemize}
  \item \textbf{Ομώνυμα σε άσχετους τομείς.} Το άρθρο \en{«FASTRAC 1»} περιγράφει
  διαστημικό δορυφόρο, ενώ το \en{«JCB Fastrac»} περιγράφει γεωργικό ελκυστήρα· τα δύο
  κείμενα μοιράζονται τη λέξη \en{«Fastrac»} αλλά ανήκουν σε εντελώς διαφορετικά
  θεματικά πεδία. Η πραγματική ετικέτα είναι θετική (πιθανότατα σύνδεσμος
  αποσαφήνισης λόγω κοινού ονόματος), αλλά η πρόβλεψη είναι αρνητική: η ομοιότητα
  ενσωματώσεων σωστά εντοπίζει τη θεματική απόσταση, λάθος όμως ως προς το αν
  υπάρχει ακμή.
  \item \textbf{Ψευδώς θετικά από κοινό πρότυπο \en{infobox}.} Δύο άρθρα ειδών
  σκαθαριού του γένους \en{Abacetus} (\en{«Abacetus laevigatus»} και ένα δεύτερο
  είδος) μοιράζονται πυκνή ταξινομική ορολογία και πανομοιότυπη δομή προτύπου
  \en{taxobox}, χωρίς καμία πραγματική ακμή μεταξύ τους. Η επιφανειακή ομοιότητα
  κειμένου, προερχόμενη από κοινό μορφότυπο και όχι από κοινό περιεχόμενο, οδηγεί σε
  ψευδώς θετική πρόβλεψη.
  \item \textbf{Χαλαρή θεματική συνάφεια.} Το άρθρο αποσαφήνισης \en{«poor me»}
  (τίτλοι τραγουδιών) συνδέεται πραγματικά με το \en{«victim mentality»} (ψυχολογική
  έννοια), μια σχέση νοήματος-προς-τίτλο που απαιτεί βαθύτερη συλλογιστική από την
  ομοιότητα ενσωματώσεων παραγράφου· η πρόβλεψη είναι αρνητική.
\end{itemize}

Στις δύο πρώτες περιπτώσεις η κατεύθυνση του σφάλματος είναι αντίθετη (ψευδώς αρνητικό
έναντι ψευδώς θετικού), αλλά η υποκείμενη αιτία είναι κοινή: η ομοιότητα κειμένου σε
επίπεδο ενσωμάτωσης δεν ταυτίζεται με την ύπαρξη ακμής σε αυτό το σύνολο δεδομένων.
Δύο άρθρα με κοινό όνομα αλλά διαφορετικό αντικείμενο μπορεί να συνδέονται (σύνδεσμος
αποσαφήνισης), ενώ δύο άρθρα με κοινή μορφοποίηση \en{infobox} αλλά ανεξάρτητο
περιεχόμενο μπορεί να μη συνδέονται καθόλου. Αυτό συνάδει με την τάση υποπρόβλεψης
θετικής ετικέτας που παρατηρήθηκε στη \S\ref{sec:cascade-results} (30{,}8\%
προβλεπόμενα θετικά έναντι 78{,}0\% πραγματικά θετικά): το Επίπεδο~3 τείνει να
ερμηνεύει τη θεματική απόσταση ως απουσία ακμής, κάτι που είναι σωστό στις
περισσότερες περιπτώσεις αλλά όχι στους συνδέσμους αποσαφήνισης/ομωνυμίας.


\section{Σύγκριση Αποτελεσματικότητας}

\subsection{Ρυθμός Επεξεργασίας}

\subsection{Ποσοστό Κλήσεων Μετασχηματιστή ανά Όριο Επιπέδου}


\section{Ανάλυση Σφαλμάτων}

Το \en{CascadeLP} κάνει 261 σφάλματα ταξινόμησης στις 187.603 ζεύγη επικύρωσης. Από
αυτά, το 93{,}49\% (244 σφάλματα) προέρχεται από ζεύγη που χαρακτηρίζονται \en{hard}
από το πρωτόκολλο χαρακτηρισμού δυσκολίας (\S\ref{sec:forensics-protocol}), και το
6{,}51\% (17 σφάλματα) από τετριμμένα ζεύγη υψηλής ομοιότητας κειμένου --- κανένα
σφάλμα δεν προέρχεται από αυτοβρόχους ή ζεύγη υψηλού \en{CN}, όπου το Επίπεδο~0 και το
Επίπεδο~1 αντίστοιχα είναι αλάνθαστα.

Το μοντέλο αναφοράς \en{TF-IDF + SVM} (\S\ref{sec:svm-baseline}), εκπαιδευμένο και
αξιολογημένο στην ίδια διαίρεση επικύρωσης, κάνει 16.696 σφάλματα --- περίπου
64 φορές περισσότερα. Το προφίλ σφαλμάτων του συγκεντρώνεται στην κατηγορία \en{hard}
σχεδόν στον ίδιο βαθμό με το \en{CascadeLP}: 93{,}63\% από ζεύγη \en{hard} και 6{,}37\%
από τετριμμένα ζεύγη υψηλής ομοιότητας κειμένου. Ο πίνακας~\ref{tab:error-comparison}
συνοψίζει τη σύγκριση.

\begin{table}[h!]
  \centering
  \begin{tabular}{lrrrr}
    \toprule
    & \multicolumn{2}{c}{\textbf{\en{CascadeLP}}} & \multicolumn{2}{c}{\textbf{\en{SVM}}} \\
    \textbf{Κατηγορία Δυσκολίας} & \textbf{Σφάλματα} & \textbf{\%} & \textbf{Σφάλματα} & \textbf{\%} \\
    \midrule
    \en{hard}                  & 244 & 93{,}49\% & 15.632 & 93{,}63\% \\
    Υψηλή Ομοιότητα Κειμένου   & 17  & 6{,}51\%  & 1.064  & 6{,}37\%  \\
    \midrule
    \textbf{Σύνολο}             & \textbf{261} & \textbf{100\%} & \textbf{16.696} & \textbf{100\%} \\
    \bottomrule
  \end{tabular}
  \caption{Σύγκριση προφίλ σφαλμάτων \en{CascadeLP} έναντι \en{SVM} ανά κατηγορία
  δυσκολίας, σύνολο επικύρωσης (n=187.603).}
  \label{tab:error-comparison}
\end{table}

Το \en{SVM} υποαποδίδει ακόμη και στην κατηγορία «τετριμμένη» υψηλής ομοιότητας
κειμένου: ακρίβεια 0{,}9838 αλλά \en{Macro F1} μόλις 0{,}4959 σε αυτό το υποσύνολο ---
το ίδιο μοτίβο κατάρρευσης λόγω ανισορροπίας κλάσεων που παρατηρήθηκε στο Επίπεδο~1
του \en{CascadeLP} (\S\ref{sec:cascade-results}). Στην κατηγορία \en{hard}, το \en{SVM}
πετυχαίνει ακρίβεια 0{,}8717 και \en{Macro F1} 0{,}7035, ενώ το \en{CascadeLP}
πετυχαίνει 0{,}9980 και 0{,}9964 αντίστοιχα στο ίδιο υποσύνολο --- μια διαφορά που δεν
εξηγείται απλά από την ομοιότητα \en{TF-IDF}, εφόσον και τα δύο μοντέλα τη
χρησιμοποιούν, αλλά από τον συνδυασμό της με τα δομικά χαρακτηριστικά του Επιπέδου~1
και ιδίως με τα χαρακτηριστικά \en{POS} του Επιπέδου~2, το οποίο απορροφά το 80{,}57\%
των ζευγών επικύρωσης με σχεδόν μηδενικό ποσοστό σφάλματος. Και τα δύο μοντέλα
συγκεντρώνουν τη μεγάλη πλειονότητα των σφαλμάτων τους στην κατηγορία \en{hard}, σε
σχεδόν πανομοιότυπη αναλογία (93{,}49\% έναντι 93{,}63\%) παρά τη 64-πλάσια διαφορά
στο απόλυτο πλήθος σφαλμάτων --- εύρημα που επιβεβαιώνει τη συνέπεια του πρωτοκόλλου
χαρακτηρισμού δυσκολίας ως προγνωστικού δείκτη πραγματικής δυσκολίας ταξινόμησης,
ανεξάρτητα από την οικογένεια μοντέλου. Ανάλυση ανά σύνολο δεδομένων αναβάλλεται
μέχρι την εκτέλεση του πρωτοκόλλου γενίκευσης (\S\ref{sec:generalization-protocol}).
